\section{Production System}\label{sec:system}

The system under evaluation is a deployed retrieval-augmented agent
serving compensation-policy questions over the Bank of Korea
Compensation Regulations. We summarize the architecture as it informs
our evaluation choices.

\subsection{Three Storage Backbones}

The same set of facts is materialized in three forms.

\paragraph{TypeDB v3 (hyper-relational graph).} The entire fact base
is loaded as a TypeQL ontology with N-ary relations and explicit
roles. A position-pay fact, for example, is stored as a single
relation \texttt{(해당직급:\$g, 해당직위:\$p, 적용기준:\$ref) isa
직책급결정}, in which the (job-grade, position, position-pay-amount)
triple co-participate in one fact rather than being decomposed into
binary edges. The schema (\texttt{schema/compensation\_regulation.tql})
defines 17 entity types, 8 N-ary relation types, and 35 attributes.

\paragraph{Neo4j (property graph).} The same data is loaded as a Neo4j
property graph in which N-ary relations are flattened: each
hyper-relational role-player tuple becomes a binary edge with extra
properties on the edge or on intermediate nodes. The Neo4j store has
9 node labels (e.g., \texttt{JobGrade}, \texttt{PositionPay},
\texttt{OverseasPay}) and 7 relationship types
(\texttt{HAS\_BASE\_SALARY}, \texttt{HAS\_POSITION\_PAY}, etc.).

\paragraph{Preprocessed context document.} A markdown
\texttt{regulation\_context.md} consolidates the article text (제2조
through 제15조), term-normalization rules (e.g., \texttt{G5 =
종합기획직원 5급}), calculation formulas, and four JSON lookup blocks
(salary differential, position pay, overseas salary, salary table).
This document is the input to the \backendname{context} backend,
which performs production-style chunked retrieval before invoking
the chat LLM.

\subsection{ReAct Agent}

The KG agents (\backendname{typedb} and \backendname{neo4j}) are
LangGraph state machines following a ReAct pattern. A typical step
proceeds:

\begin{enumerate}
\item The chat model receives the user question and a schema
      summary, and decides whether a database query is needed.
\item If a query is needed, control passes to the query model, which
      authors a TypeQL or Cypher query against the loaded schema.
\item The query is executed; results return as structured rows.
\item The chat model verbalizes the results into a natural-language
      answer, optionally citing schema rules.
\end{enumerate}

A failed parse triggers retry with the parse error fed back into the
prompt; the agent budget is capped at three query attempts per
question.

\subsection{Decoupled Query Agent (Setting B)}

The production code separates the chat model
(\texttt{create\_chat\_model()}) from the query model
(\texttt{create\_qwen\_model()}) by design: the original deployment
uses Qwen-Coder for query authoring even when the chat model is HCX.
For our evaluation we honour this separation by holding the query
model fixed at \texttt{Qwen3-Coder-30B-A3B-Instruct} and varying only
the chat model.

This decoupling matters: a pilot in which the same LLM played both
roles produced highly asymmetric results across LLM families,
confounding the KG vs.\ context comparison. The pilot is reported as
Setting A in Appendix~\ref{sec:appendix}; all main numbers in this
paper are Setting B.

\subsection{In-Context Surface-Form Tiers}

To isolate the contribution of the query-generation stage from the
intrinsic readability of the KG, we add four in-context backends that
side-step the agent entirely. Each constructs a single fixed prompt
of the form

\begin{quote}\small
\texttt{[article-text prefix common to all tiers]} \\[2pt]
\texttt{=== 자료 시작 ===} \\
\texttt{[KG dump in tier-specific surface form]} \\
\texttt{=== 자료 끝 ===}
\end{quote}

and invokes the chat model with this prompt as the system message and
the user question as the user message. The four surface forms (TQL,
Cypher, NL, JSON-LD) are produced from the same underlying source
tables (\texttt{eval/dump\_context\_tiers.py}), so any difference in
accuracy between tiers is purely a surface-form effect.
